\section{Injectivity-length reduction for the closure property}\label{sec:block_word_stripping}

This section implements the inverting-and-growing-back reduction from
\cite[Section~IV.C]{Cirac2021Matrix}: a commutation or annihilation relation
for length-$m$ word products is reduced to the injectivity length $L_0$, then
used in the periodic-boundary closure-property step. Throughout,
$A^w := A^{i_1}\cdots
A^{i_L}$ denotes the product along a word $w = (i_1,\ldots,i_L)$, $\Gamma_N$
sends a boundary matrix to its open-chain vector, and $A$ being
\emph{$L_0$-block-injective} means $\spn\{A^u : |u|=L_0\} = \MN{D}$, so a
relation valid on every length-$L_0$ product extends by linearity to all of
$\MN{D}$.

\begin{theorem}[Common right factor from suffix equations]
    \label{thm:block_word_strip_right_factor}
    \lean{MPSTensor.exists_right_factor_of_block_word_compatibility}
    \leanok
    \uses{def:l_blk_injective}
    Assume $A$ is $L_0$-block-injective with $L_0 > 0$. Let $\{Z_u\}_{|u|=L_0}$
    be a family indexed by words of length $L_0$. If for some $K \ge 0$ and every
    suffix word $w$ of length $K$ there exists $Y_w \in \MN{D}$ such that,
    for every word $u$ of length~$L_0$,
    \[
        Z_u A^w = A^u Y_w,
    \]
    then there exists $X \in \MN{D}$ with $Z_u = A^u X$ for every word $u$ of
    length $L_0$.
\end{theorem}

\begin{proof}
    \leanok
    Induct on $K$. For $K = 0$ the empty suffix gives $Z_u = A^u Y_{\emptyset}$.
    For the successor step, fix the first physical index $i$ and apply the
    induction hypothesis to the shortened suffix. This gives a matrix $X_i$
    satisfying
    \[
        Z_u A^i = A^u X_i
    \]
    for every word $u$ of length $L_0$. Extending this letter equation to all
    length-$L_0$ words and decomposing $\Id$ in the spanning family
    $\{A^u : |u| = L_0\}$ yields a common factor $X$.
\end{proof}

\begin{theorem}[Common boundary matrix from spanning word identities]
    \label{thm:word_span_common_boundary_matrix}
    \lean{MPSTensor.exists_common_boundary_matrix_of_word_identities_of_wordSpan_eq_top}
    \leanok
    \uses{def:word_span}
    Let $A$ be a tensor whose length-$K$ word products span $\MN{D}$. Let
    $\{F_b\}_{b\in B}$ and $\{Z_b\}_{b\in B}$ be two families of matrices. If,
    for every length-$K$ word $w$, there is a matrix $Y_w$ such that, for every
    $b\in B$,
    \[
        Z_b A^w=F_bY_w,
    \]
    then there is a single matrix $Y$ such that, for every $b\in B$,
    \[
        Z_b=F_bY.
    \]
\end{theorem}

\begin{proof}
    \leanok
    Suppose $Z_bM_1=F_bY_1$ and $Z_bM_2=F_bY_2$ for every $b\in B$. Then
    \[
        Z_b(M_1+M_2)
        =
        Z_bM_1+Z_bM_2
        =
        F_b(Y_1+Y_2).
    \]
    For a scalar $c$,
    \[
        Z_b(cM_1)
        =
        cZ_bM_1
        =
        F_b(cY_1).
    \]
    Thus the displayed identity holds for every matrix in the span of the
    length-$K$ word products. Since this span is all of $\MN{D}$, apply it
    to the identity matrix.
\end{proof}

\begin{theorem}[Injectivity-length commutation from length-$m$ commutation]
    \label{thm:block_word_strip_commutes}
    \lean{MPSTensor.commutes_block_words_of_commutes_long_words_of_isNBlkInjective}
    \leanok
    \uses{thm:block_word_strip_right_factor, def:l_blk_injective}
    If $A$ is $L_0$-block-injective with $L_0 > 0$, $m \ge L_0$, and
    $X \in \MN{D}$ commutes with every length-$m$ product $A^\omega$, then $X$
    already commutes with every length-$L_0$ product $A^u$.
\end{theorem}

\begin{proof}
    \leanok
    Write each length-$m$ word as a concatenation $uw$ with $|u| = L_0$ and
    $|w| = m - L_0$. The hypothesis gives
    \[
        X A^u A^w = A^u A^w X = A^u (A^w X).
    \]
    Thus the family $Z_u := X A^u$ satisfies the hypothesis of
    Theorem~\ref{thm:block_word_strip_right_factor}. One obtains
    $X A^u = A^u R$ for all $|u| = L_0$. Decomposing $\Id$ in the span of the
    length-$L_0$ words shows $R = X$, hence $X A^u = A^u X$.
\end{proof}

\begin{theorem}[Centrality from length-$m$ commutation]
    \label{thm:block_word_strip_central}
    \lean{MPSTensor.commutes_all_of_commutes_long_words_of_isNBlkInjective}
    \leanok
    \uses{thm:block_word_strip_commutes, def:l_blk_injective}
    If $A$ is $L_0$-block-injective with $L_0 > 0$, $m \ge L_0$, and
    $X \in \MN{D}$ commutes with every length-$m$ word product, then $X$
    commutes with every matrix in $\MN{D}$.
\end{theorem}

\begin{proof}
    \leanok
    By Theorem~\ref{thm:block_word_strip_commutes}, $X$ commutes with every
    length-$L_0$ word product. These products span $\MN{D}$ by block
    injectivity, so linearity gives $XM = MX$ for every $M \in \MN{D}$.
\end{proof}

\begin{lemma}[Amplifying fixed-length word commutation]
    \label{lem:word_commutation_multiple_lengths}
    \lean{MPSTensor.commutes_words_mul_of_commutes_words}
    \leanok
    If a boundary matrix $X$ commutes with every word product of length $m$, then
    it commutes with every word product whose length is a multiple $q\,m$.
\end{lemma}

\begin{proof}
    \leanok
    Proceed by induction on $q$. For $q=0$, the claim is trivial. For the
    inductive step, let a word of length $(q+1)m$ be given and write it as a
    length-$m$ prefix followed by a length-$q\,m$ suffix. The prefix commutes with
    $X$ by hypothesis, and the suffix commutes by the induction hypothesis.
\end{proof}

\begin{lemma}[Padded complement annihilation]
    \label{lem:padded_complement_annihilation}
    \lean{MPSTensor.eq_zero_of_mul_evalWord_eq_zero_of_isNBlkInjective_of_le_mul,
        MPSTensor.eq_zero_of_evalWord_mul_eq_zero_of_isNBlkInjective_of_le_mul}
    \leanok
    \uses{def:l_blk_injective}
    If $A$ is $L_0$-block-injective, $q \ge 1$, and $k \le qL_0$, then the
    two one-sided annihilation statements are:
    \[
        Z A^w = 0 \quad (|w|=k) \qquad\Longrightarrow\qquad Z=0.
    \]
    The left-handed form is
    \[
        A^w Z = 0 \quad (|w|=k) \qquad\Longrightarrow\qquad Z=0.
    \]
\end{lemma}

\begin{proof}
    \leanok
    The hypothesis $\spn\{A^u: |u|=L_0\}=\MN{D}$ implies the same spanning
    statement at every positive multiple $qL_0$. Each length-$qL_0$ word factors
    as a length-$k$ prefix followed by padding, so $Z$ annihilates all generators
    of the full span and hence annihilates the identity.
    For the left-handed form, factor each length-$qL_0$ word as a prefix followed
    by a length-$k$ suffix and use the same spanning argument.
\end{proof}

\begin{lemma}[One-sided boundary matrices are unique]
    \label{lem:one_sided_boundary_witness_unique}
    \lean{MPSTensor.right_witness_unique_of_isNBlkInjective}
    \lean{MPSTensor.left_witness_unique_of_isNBlkInjective}
    \leanok
    \uses{def:l_blk_injective, lem:padded_complement_annihilation,
        lem:wordspan_blk_inj}
    Let $A$ be $L_0$-block-injective with $L_0>0$.  If two boundary matrices
    have the same product with every one-site tensor on one fixed side, then the
    two boundary matrices are equal.
\end{lemma}

\begin{proof}
    \leanok
    If $Y_1 A^j = Y_2 A^j$ for every physical index $j$, then for every
    nonempty word $w$, induction on the word gives
    \[
        Y_1 A^w = Y_2 A^w.
    \]
    Since $A$ is $L_0$-block-injective,
    \[
        \spn\{A^w : |w|=L_0\}=\MN{D}.
    \]
    Thus the left multiplication maps $M\mapsto Y_1 M$ and
    $M\mapsto Y_2 M$ are equal on $\MN{D}$, and evaluating on the identity
    gives $Y_1=Y_2$.  The other one-sided statement is the same argument, with
    right multiplication in place of left multiplication.
\end{proof}

\begin{lemma}[Generator commutation from length-$m$ commutation]
    \label{lem:boundary_matrix_commutes_long_words}
    \lean{MPSTensor.boundary_matrix_commutes_of_isNBlkInjective_of_long_word_commutes}
    \leanok
    \uses{thm:block_word_strip_central}
    Under the hypotheses of Theorem~\ref{thm:block_word_strip_central}, one has
    $X A^i = A^i X$ for every physical index $i$.
\end{lemma}

\begin{proof}
    \leanok
    Apply Theorem~\ref{thm:block_word_strip_central} with $M = A^i$.
\end{proof}

\begin{theorem}[MPS-line containment from length-$m$ commutation]
    \label{thm:ground_space_map_mpv_of_long_word_commutes}
    \lean{MPSTensor.groundSpaceMap_mem_mpvSubmodule_of_isNBlkInjective_of_long_word_commutes}
    \leanok
    \uses{def:mpv_submodule, thm:block_word_strip_central}
    Let $A$ be injective after blocking $L_0>0$ sites, and let $m \ge L_0$.
    If a boundary matrix $X$ commutes with every length-$m$ word product, then
    \[
        \Gamma_N(X)\in\spn\{V^{(N)}(A)\}.
    \]
\end{theorem}

\begin{proof}
    \leanok
    Theorem~\ref{thm:block_word_strip_central} makes $X$ commute with the whole
    matrix algebra. The center of $\MN{D}$ consists of scalar matrices, so
    $X=c\Id$, and therefore $\Gamma_N(X)=cV^{(N)}(A)$.
\end{proof}

\begin{theorem}[MPS-line containment from positive-length commutation]
    \label{thm:ground_space_map_mpv_of_positive_word_commutes}
    \lean{MPSTensor.groundSpaceMap_mem_mpvSubmodule_of_isNBlkInjective_of_positive_word_commutes}
    \leanok
    \uses{lem:word_commutation_multiple_lengths,
        thm:ground_space_map_mpv_of_long_word_commutes}
    Let $A$ be injective after blocking $L_0>0$ sites. If $X$ commutes with
    every word product of some positive length $m$, then
    \[
        \Gamma_N(X)\in\spn\{V^{(N)}(A)\}.
    \]
\end{theorem}

\begin{proof}
    \leanok
    Lemma~\ref{lem:word_commutation_multiple_lengths} amplifies the length-$m$
    commutation hypothesis to length $L_0\,m \ge L_0$. Then
    Theorem~\ref{thm:ground_space_map_mpv_of_long_word_commutes} applies.
\end{proof}

\begin{theorem}[MPS-line containment from two one-sided equations]
    \label{thm:ground_space_map_mpv_of_common_middle}
    \lean{MPSTensor.groundSpaceMap_mem_mpvSubmodule_of_isNBlkInjective_of_two_sided_middle_compatibility}
    \leanok
    \uses{lem:common_middle_same_witness_commutation,
        thm:ground_space_map_mpv_of_positive_word_commutes}
    Let $A$ be injective after blocking $L_0>0$ sites. Suppose there are matrices
    $Y_\mu$ such that
    \[
        A^\mu A^b X = Y_\mu A^b,
        \qquad
        X A^a A^\mu = A^a Y_\mu,
    \]
    for every pair of physical letters $a,b$. Then
    \[
        \Gamma_N(X)\in\spn\{V^{(N)}(A)\}.
    \]
\end{theorem}

\begin{proof}
    \leanok
    The identities in~(\ref{eq:ph_common_middle_one_sided}) give commutation
    with all words of length $|\mu|+2$, a positive length. The positive-length
    containment theorem then applies.
\end{proof}

\begin{theorem}[Boundary equality implies containment in the MPS line]
    \label{thm:ground_space_map_mpv_of_wrapped_witness_comparison}
    \lean{MPSTensor.groundSpaceMap_mem_mpvSubmodule_of_isNBlkInjective_of_wrapped_witness_comparison}
    \leanok
    \uses{lem:two_sided_middle_of_wrapped_witness_comparison,
        thm:ground_space_map_mpv_of_common_middle}
    Let $A$ be injective after blocking $L_0>0$ sites, and fix a physical letter
    $\eta$ used outside the complementary sites. Suppose that, for every
    boundary condition $\tau$ and every physical letter $j$,
    \[
        C^+_\tau A^jX=Y^+_\tau A^j,
        \qquad
        XA^jC^-_\tau=A^jY^-_\tau,
    \]
    where $C^+_\tau$ and $C^-_\tau$ are the complementary words exposed by the
    two boundary-crossing cyclic windows. Suppose also that the two boundary
    matrices satisfy
    \[
        Y^+_{\tau^+_\eta(\mu)} = Y^-_{\tau^-_\eta(\mu)},
    \]
    for every word $\mu$ on the complementary sites, then
    \[
        \Gamma_N(X)\in\spn\{V^{(N)}(A)\}.
    \]
\end{theorem}

\begin{proof}
    \leanok
    Lemma~\ref{lem:two_sided_middle_of_wrapped_witness_comparison} gives a
    family $Y_\mu$ satisfying
    \[
        A^\mu A^bX=Y_\mu A^b,
        \qquad
        XA^aA^\mu=A^aY_\mu.
    \]
    Theorem~\ref{thm:ground_space_map_mpv_of_common_middle} applies.
\end{proof}

\begin{theorem}[Closure-property step for periodic chains]
    \label{thm:conditional_normal_range_reduction_witness_comparison}
    \lean{MPSTensor.chainGroundSpace_le_mpvSubmodule_of_isNBlkInjective_of_wrapped_witness_comparison}
    \leanok
    \uses{thm:cyclic_normal_open_chain_range_reduction,
        thm:reduced_wrapped_boundary_compatibilities,
        thm:ground_space_map_mpv_of_wrapped_witness_comparison}
    Let $A$ be injective after blocking $L_0>0$ sites and let $D \ge 1$. Assume
    $N \ge 2$ and $L_0 < L \le N$. Suppose that, for every physical letter
    $\eta$ used in the two boundary conditions and every boundary matrix
    obtained from an $N$-site chain ground state, the identities
    \[
        A^\mu A^j X = Y^+_{\tau^+_\eta(\mu)}A^j,
        \qquad
        X A^j A^\mu = A^jY^-_{\tau^-_\eta(\mu)},
    \]
    hold for every physical letter $j$ and every complementary word $\mu$, and
    are supplemented, for every $\mu$, by
    \[
        Y^+_{\tau^+_\eta(\mu)}
        =
        Y^-_{\tau^-_\eta(\mu)}.
    \]
    Then
    \[
        \mathcal G_{N,L}(A)\subseteq \spn\{V^{(N)}(A)\}.
    \]
\end{theorem}

\begin{proof}
    \leanok
    \uses{thm:cyclic_normal_open_chain_range_reduction,
        thm:reduced_wrapped_boundary_compatibilities,
        thm:ground_space_map_mpv_of_wrapped_witness_comparison}
    The cyclic-to-open-chain reduction writes any chain ground state as
    $\Gamma_N(X)$. The two boundary-crossing cyclic windows give
    \[
        A^\mu A^jX=Y^+_{\tau^+_\eta(\mu)}A^j,
        \qquad
        XA^jA^\mu=A^jY^-_{\tau^-_\eta(\mu)}.
    \]
    Together with
    \[
        Y^+_{\tau^+_\eta(\mu)}
        =
        Y^-_{\tau^-_\eta(\mu)},
    \]
    Theorem~\ref{thm:ground_space_map_mpv_of_wrapped_witness_comparison}
    finishes.
\end{proof}
